\begin{abstract}
Key-Value (KV) cache reduces inference latency in large language models (LLMs). Traditional prefix-based reuse has low cache hit rates across inference requests because it requires exact token and position matches. To improve efficiency, recent system optimizations introduce position-independent KV reuse, allowing KV cache to be reused whenever identical text chunks appear, regardless of their position in the sequence. 

We show this design introduces a new threat, \emph{KV Cache Hijacking}. Since KV caches are retrieved by token match but encode the context in which they were originally computed, the KV tied to a benign-looking token chunk may encode an attacker-controlled prefix. When later reused in a victim query, this contaminated KV silently hijacks the model's behavior, even if \emph{no attacker-controlled text appears in the input.}

We introduce \sys, the \emph{first} attack framework that systematically exploits this vulnerability, demonstrating its severity and practicality. \sys optimizes an attacker-controlled prefix, so that the KV computed for a subsequent common benign text encodes the attacker's goal, while the text remains unchanged for future cache hits. \sys achieves an average 94\% success rate in a single attempt, remains effective under realistic constraints including low hit rates (10\%) and frequent recomputation (50\%), persists over multi-turn interactions, and transfers across models in black-box settings. We further provide design insights for building secure KV reuse systems.\footnote{The source code of \sys is publicly available: \url{https://github.com/YichiCS/KV-Cache-Hijack}} 
\end{abstract}